% SGA 6, Exposición XIV: traducción española alineada con la fuente.
% Alcance exacto: fr.tex líneas 15447--15554, secciones 1--3 completas.
% Las lecturas fuente corregidas se conservan mediante \sisi.

\providecommand{\sisi}[2]{#2}

\section*{Exposición XIV. Problemas abiertos en teoría de intersecciones}
\begin{center}
por A. Grothendieck
\end{center}

En la presente exposición pasamos revista a algunos de los problemas
relacionados con las cuestiones encontradas en este seminario y que
merecerían ser estudiados. Dejamos de lado los problemas relacionados con la
\emph{teoría local} de intersecciones de Serre [14], que no hemos tenido que
utilizar en el Seminario, y remitimos especialmente a loc. cit. para las dos
conjeturas fundamentales que permanecen abiertas en esta dirección
(loc. cit., cap. V, \S~4). Omitimos también los problemas, muy importantes
para las aplicaciones aritméticas, relacionados con la \emph{estructura de
los grupos de clases de ciclos}, especialmente sobre un esquema algebraico
proyectivo y liso; para su discusión remitimos a S. Kleiman,
\emph{Algebraic cycles and the Weil conjectures}, en \emph{Dix exposés sur
la cohomologie des schémas} (Masson--North Holland Publishing Cie). En el
número 4 de esta exposición damos asimismo algunos complementos al
Seminario, relativos en particular al \emph{anillo de Chow}, que no habían
encontrado su lugar en las exposiciones precedentes.

\section*{1. Operaciones $\Lambda^i$ en la categoría derivada
$D(X)$\footnote{Cf. la nota de la página 23.}}

Sea $X$ un topos localmente anillado y sea $K^\bullet(X)$ el anillo de las
clases de complejos perfectos sobre $X$ (Exp. IV). Se querría definir en
$K^\bullet(X)$ una estructura de $\lambda$-anillo (Exp. V), funtorial en
$X$, tal que el homomorfismo natural
\[
K^\bullet_{\mathrm{naif}}(X)\longrightarrow K^\bullet(X)
\]
(donde el primer miembro es el anillo de las clases de módulos localmente
libres sobre $X$, provisto de la $\lambda$-estructura definida en la Exp. V)
sea un $\lambda$-homomorfismo. Con este fin, habría que definir funtores, no
aditivos en general, de «potencia exterior»
\[
\Lambda^i:D^-(X)\longrightarrow D^-(X),
\]
que transformen complejos pseudocoherentes (resp. perfectos) en complejos
pseudocoherentes (resp. perfectos) y posean, respecto de los triángulos
exactos, las propiedades formales necesarias para definir una
$\lambda$-estructura sobre $K^\bullet(X)$ mediante la fórmula
\begin{equation}\tag{1.1}
\lambda^i(\cl^\bullet(L^\bullet))=\cl^\bullet(\Lambda^i L^\bullet).
\end{equation}

Cuando $X$ es de característica cero, es decir, cuando su haz estructural es
un haz de $\Q$-álgebras, el problema se resuelve sencillamente tomando
\begin{equation}\tag{1.2}
\Lambda^i L^\bullet=\bigl(\overset{i}{\otimes}{}^{\mathbb L}L^\bullet\bigr)^{\mathrm{alt}},
\end{equation}
donde la potencia tensorial $i$-ésima, en el sentido de las categorías
derivadas, de $L^\bullet$ se considera como un objeto de
$D^-(\mathcal O_i)$, siendo $\mathcal O_i$ el álgebra del grupo simétrico
$\mathfrak S_i$ con coeficientes en $\mathcal O_X$, y donde el exponente
alt denota el paso a la «parte alternada bajo la acción de
$\mathfrak S_i$» (operación que tiene sentido gracias a la hipótesis de
característica cero). Una fórmula de este tipo ya se utilizó en [8], citado
en lo sucesivo como [RRR], para la primera demostración del teorema de
Riemann--Roch.

Sin hipótesis sobre la característica, cuando $L^\bullet$ es nulo en grados
$>0$, Dold--Puppe [6] muestra que $L^\bullet$ es canónicamente homotópico al
complejo asociado a un módulo semisimplicial $L^{\prime\bullet}$ asociado
canónicamente a $L^\bullet$. Un candidato bastante natural para
$\Lambda^iL^\bullet$, digamos para $L^\bullet$ plano, sería entonces el
complejo asociado al módulo semisimplicial $\Lambda^iL^{\prime\bullet}$,
cuya componente de grado $n$ es $\Lambda^iL'_n$. Pero, para extender la
definición a complejos de $D^-(X)$ que ya no se suponen nulos en grados
$>0$, habría que estudiar el comportamiento de la operación anterior
respecto de la traslación $L^\bullet\mapsto L^\bullet[1]$ (o
\emph{suspensión}, en la terminología de [6]); desgraciadamente, no parece
que se disponga de un isomorfismo
\[
\Lambda^i(L^\bullet[1])\simeq \Lambda^i(L^\bullet)[i]
\]
que permitiese efectuar trivialmente la extensión deseada.

\section*{2. La fórmula de Riemann--Roch sin hipótesis proyectivas}

\textbf{2.1.} Una vez en posesión de una $\lambda$-estructura natural sobre
los $K^\bullet(X)$, es posible formular el teorema de
Riemann--Roch, en la forma considerada en la Exp. VIII, para cualquier
morfismo propio de intersección completa $f:X\to Y$, e incluso, más
generalmente, para un esquema relativo $X$ sobre un topos localmente
anillado $Y$, sometiendo todavía $X$ a las condiciones relativas
correspondientes. La clase de Todd de $f$ se define mediante la clase en
$K^\bullet(X)$ del complejo cotangente relativo $L^\bullet$ de $f$,
construido en la Exp. 0, 4.4; este complejo es aquí perfecto gracias a la
hipótesis local hecha sobre $f$. Parece claro que la demostración de la
fórmula de Riemann--Roch en este caso general exigirá la introducción de
ideas esencialmente nuevas, y es posible que estas mismas ideas proporcionen
la clave de una demostración de la fórmula análoga en el caso analítico
complejo (¡o analítico rígido!). Recordemos que, en el caso analítico
complejo, solo se conoce actualmente el caso de un haz localmente libre $F$
sobre una variedad no singular compacta $X$, sobre un $Y$ puntual, gracias
a la fórmula del índice de Atiyah--Singer; el caso en que $F$ solo se supone
coherente parece escapar todavía a las técnicas conocidas.

\textbf{2.2.} Como ingrediente esencial de la fórmula de Riemann--Roch,
recordemos también el \emph{problema de finitud} planteado en la Exp. III:
si $f:X\to Y$ es un morfismo propio y $L^\bullet$ un complejo sobre $X$
pseudocoherente relativamente a $Y$, ¿es $Rf_*(L^\bullet)$
pseudocoherente? Una demostración general de esta conjetura pondrá
probablemente en juego ideas algebraicas que permitirán demostrar el
«teorema» de finitud de Grauert [7] bajo condiciones sensiblemente más
generales que las de loc. cit., situándose desde el principio sobre espacios
o topoi de base topológicamente anillados en $\C$-álgebras $Y$, más generales
que los espacios analíticos. Esto abarcaría el caso en que $Y$ fuese un
espacio localmente compacto con el haz de funciones continuas complejas, una
variedad $C^\infty$ con el haz de funciones complejas $C^\infty$, un espacio
analítico real, etc.

\textbf{2.3.} Señalemos, finalmente, que en estas cuestiones convendría
reemplazar la hipótesis de que $f$ sea propio por la hipótesis de que los
haces de cohomología $H^i(L^\bullet)$ de los complejos considerados sobre
$X$ tengan \sisi{sus grupos propios}{soportes propios} sobre $Y$.

\section*{3. Fórmula de Riemann--Roch «sin denominadores» para una inmersión}

\textbf{3.1.} Sea $f:Y\to X$ una inmersión cerrada regular (Exp. VII), con
haz conormal $N$. Supongamos, para simplificar, que $X$ admite un haz
invertible amplio. La fórmula de Riemann--Roch para $y\in K^\bullet(Y)$
\[
\ch_X(f_*(y))=f_*\bigl(\ch_Y(y)\operatorname{Todd}(-\check N)\bigr),
\]
junto con la fórmula general (Exp. VII)
\[
f_*(y')\,f_*(y'')=f_*\bigl(y'y''c_d(\underline N)\bigr),\qquad
y',y''\in\operatorname{Gr}^\bullet(Y)_\Q,
\]
suponiendo que $f$ tiene codimensión $d$, permite expresar las clases de
Chern de $f_*(y)$ en $\operatorname{Gr}^\bullet(X)_\Q$ en la forma
\begin{equation}\tag{3.1}
c_i(f_*(y))=f_*\bigl(P_i((c_\alpha(y))_\alpha,
(c_\beta(N))_\beta)\bigr),
\end{equation}
donde $P_i$ es un polinomio universal en los $c_\alpha(y)$ y $c_\beta(N)$,
a priori con coeficientes racionales. En realidad, se comprueba en RRR que
estos polinomios tienen coeficientes \emph{enteros}. Esto plantea la cuestión
de la validez de una fórmula de Riemann--Roch «sin denominadores».

En general, una fórmula semejante para $c_i(f_*(y))$ no puede buscarse en el
propio $\operatorname{Gr}^\bullet(X)$, puesto que el homomorfismo
$f_*:K^\bullet(Y)\to K^\bullet(X)$ no es compatible con las
$\lambda$-filtraciones, ni siquiera tras desplazarlas en $d$ (4.6), sino
solo salvo torsión. Por ello no es posible definir un homomorfismo
\[
f_*:\operatorname{Gr}^\bullet(Y)\longrightarrow
\operatorname{Gr}^\bullet(X)
\]
que induzca el homomorfismo natural
$\sisi{f}{f_*}:\operatorname{Gr}^\bullet(Y)_\Q\to
\operatorname{Gr}^\bullet(X)_\Q$ y permita interpretar (3.1) como una
fórmula en $\operatorname{Gr}^\bullet(X)$, y no solo en
$\operatorname{Gr}^\bullet(X)_\Q$.

Pero, si $X$ es regular, de modo que
$K^\bullet(X)\simeq K_\bullet(X)$ (IV 2.5), se puede dotar
$K^\bullet(X)$ de la filtración topológica por la codimensión de los
soportes, heredada de la filtración correspondiente de $K_\bullet(X)$. Es
plausible que esta filtración sea compatible con la estructura de anillo;
esto puede demostrarse, al menos en RRR, cuando $X$ es liso sobre un cuerpo
$k$, utilizando la teoría del anillo de Chow [4]. Cuando esto ocurre, el
graduado asociado
\[
\Gtop^\bullet(X)
\]
es un \emph{anillo graduado}; suponiendo además que $Y$ es regular, se dispone
entonces de un homomorfismo natural de grado $d$ de grupos graduados
\[
f_*=f_*^{\mathrm{top}}:\Gtop^\bullet(Y)\longrightarrow
\Gtop^\bullet(X).
\]
Como la $\lambda$-filtración de $K^\bullet(X)$ es más fina que la filtración
topológica por codimensión que acabamos de considerar (X 1.3.3), se obtiene
un homomorfismo canónico de grupos graduados, e incluso de anillos graduados,
\[
\operatorname{Gr}^\bullet(X)\longrightarrow\Gtop^\bullet(X),
\]
que es, por lo demás, un isomorfismo módulo torsión (Exp. VIII). Esto permite,
como en RRR, considerar sobre $X$ una teoría de clases de Chern con valores
en $\Gtop^\bullet(X)$, y análogamente sobre $Y$, dando así sentido a (3.1)
como fórmula en $\Gtop^\bullet(X)$. En esta forma, la fórmula se demuestra
en RRR cuando $X$ es de \emph{característica cero}, al menos si $X$ es liso
sobre un cuerpo; esta última hipótesis solo sirve para asegurar que la
filtración topológica de $K^\bullet(X)$ sea compatible con la estructura de
anillo. La hipótesis de característica cero se utiliza esencialmente para
disponer de fórmulas del tipo (1.1), (1.2), y la misma demostración debería
aplicarse esencialmente sin hipótesis sobre la característica una vez
resuelta afirmativamente la cuestión planteada en el número 1.

\textbf{3.2.} Cuando $X$ es liso sobre un cuerpo, se dispone también del
anillo de Chow
\[
A(X),
\]
y de una teoría de clases de Chern sobre $X$ con valores en $A(X)$ [9], más
fina que la teoría con valores en $\Gtop^\bullet(X)$ recién introducida; las
últimas clases de Chern son imágenes de las clases de Chern de estilo Chow
por el homomorfismo canónico sobreyectivo (cf. 4.1)
\[
\varphi:A(X)\longrightarrow\Gtop^\bullet(X).
\]
Este homomorfismo es un isomorfismo módulo torsión, pero no un isomorfismo en
general (4.7). Dicho esto, si $X$ e $Y$ son lisos, la inclusión $f:Y\to X$
define todavía un homomorfismo de grado $d$ de grupos graduados
\[
f_*:A^\bullet(Y)\longrightarrow A^{\bullet+d}(X),
\]
lo que permite interpretar (3.1) como una fórmula en $A^\bullet(X)$. Esta
fórmula se reduce, por lo demás, a $0=0$ para $i<d$, y por tanto solo
presenta interés para $i\sisi{>}{\ge}d$. Incluso para $i=d=2$, la fórmula no está % Canon SGA-SRC-001404
demostrada actualmente (cf. 4.3), como tampoco la fórmula menos precisa que
se deduce de ella reemplazando la equivalencia racional de los ciclos
algebraicos por la equivalencia algebraica. A diferencia de la fórmula
análoga en $\Gtop^\bullet(X)$, no parece haber una razón convincente para
suponer que la fórmula sin denominadores deba ser cierta en $A(X)$. Una
buena manera de probarla, dada nuestra ignorancia casi total de la naturaleza
de los grupos $A^i(X)$ para $i>2$, parece consistir en probar las fórmulas
que pueden deducirse en grupos $A^1(Z)=\Pic(Z)$, por ejemplo transformando
los dos miembros de (3.1) mediante
\[
g_*:A(X)\longrightarrow A(Z),
\]
donde $g:X\to Z$ es un morfismo propio sobre un $Z$ cuasiproyectivo liso de
dimensión $\dim X-i+1$.

\emph{Nota añadida.} La fórmula, así como la variante señalada en 3.3, fue
demostrada posteriormente por P. JOUANOLOU (noviembre de 1969).

\textbf{3.3.} Finalmente, para dar sentido a (3.1), se puede utilizar también
la teoría cohomológica de clases de Chern (SGA 5 VII), siempre que se disponga
de \emph{homomorfismos de Gysin}
\[
f_*:H^p(Y,\mu_n^{\otimes i})\longrightarrow
H^{p+2d}(X,\mu_n^{\otimes(i+d)}).
\]
(Cf. n\textsuperscript{o}~6). Este es el caso, al menos, si $Y$ y $X$ son
ambos lisos sobre un esquema $S$, gracias al teorema de pureza cohomológica
relativa (SGA 4 XVI 3.7), y, conjeturalmente, si $Y$ y $X$ son ambos
regulares, según el teorema de pureza cohomológica absoluta, demostrado al
menos para $X$ excelente de característica cero (SGA 4 XIX 3.2). Cuando $X$
e $Y$ son lisos sobre el cuerpo de los complejos, se puede utilizar también
la clase de Chern entera sobre el espacio localmente compacto $X(\C)$ de los
puntos complejos de $X$. La fórmula (3.1) es establecida en este caso por
\uline{Atiyah} y \uline{Hirzebruch} [1], quienes tratan más generalmente una
inmersión de variedades analíticas complejas, trabajando con el $K'$
«ingenuo». De ello resulta que la fórmula $\ell$-ádica (3.1) es válida
siempre que $Y$ y $X$ sean lisos sobre un cuerpo \uline{algebraicamente
cerrado y de característica cero}, utilizando el «principio de Lefschetz» y
los resultados generales de SGA 4. Esto hace extremadamente plausible la
validez de la fórmula cohomológica (3.1); obsérvese, sin embargo, que no está
demostrada actualmente, ni siquiera si $X$ e $Y$ son lisos sobre un cuerpo
$k$ de característica cero, cuando $k$ no es algebraicamente cerrado.

\textbf{3.4.} También se puede trabajar con la teoría cohomológica de clases
de Chern de estilo De Rham o Hodge [13]. El problema de la validez de (3.1)
en este caso solo se plantea si $X$ no es de característica cero, pues se
sabe de todos modos (n\textsuperscript{o}~6) que la fórmula en cuestión es
cierta módulo torsión. Para los esquemas lisos cuasiproyectivos sobre un
cuerpo $k$, la versión De Rham o Hodge de (3.1) sería, por lo demás,
consecuencia de la versión de (3.1) para clases de ciclos módulo equivalencia
algebraica. Así, la versión cohomológica De Rham o Hodge proporciona otra
manera accesible de probar la validez de la versión «clases de ciclos» de
(3.1).

\textbf{3.5.} La demostración [1] de la versión cohomológica trascendente de
(3.1) es de naturaleza esencialmente trascendente y utiliza entornos
tubulares de $Y$ en $X$. Esto sugiere que la demostración de una fórmula del
tipo (3.1) podría estar ligada a la consideración sistemática de grupos $K'$
y grupos de cohomología \uline{con soporte en $Y$}; dicha consideración
podría proporcionar un medio técnico suficiente de «localización alrededor
de $Y$».
